\begin{tikzpicture}
\node[note,anchor=west] at (-1.6,4.8) {历史 A：当前分配实际取回 100、101、102};
\node[good,text width=2.4cm] at (0,3.6) {100\\本地处理完成};
\node[warn,text width=2.4cm] at (3.8,3.6) {101\\本地事务阻塞};
\node[good,text width=2.4cm] at (7.6,3.6) {102\\本地处理完成};
\draw[arr] (-1.3,2.3)--node[note,below]{完成前缀止于 100，提交 next=101}(1.8,2.3);
\draw[Gold,densely dashed,thick] (1.9,4.35)--(1.9,2.05);
\node[note,anchor=west] at (-1.6,1.05) {历史 B：压缩后实际取回 100、102，读取位置为 103};
\node[good,text width=2.4cm] at (0,-.2) {100\\本地处理完成};
\node[note,text width=2.4cm] at (3.8,-.2) {101\\无待处理应用记录};
\node[good,text width=2.4cm] at (7.6,-.2) {102\\本地处理完成};
\draw[arr] (-1.3,-1.5)--node[note,below]{实际取得的记录全部完成，提交 next=103}(9,-1.5);
\draw[Blue,densely dashed,thick] (9.1,.6)--(9.1,-1.75);
\node[note,text width=11.7cm] at (3.8,-3) {协调器保存当前分配、取回记录和逐项完成状态\\101 的数值空隙由读取结果解释；已取回但未完成的 101 保留处理责任};
\end{tikzpicture}
